\section{1. Introduction}

\subsection{1.1 Motivation}

Attention dominates the compute and memory cost of transformer inference:

\[
\text{Attn}(Q, K, V) = \text{softmax}\!\left(\frac{Q K^\top}{\sqrt{d_h}}\right) V
\tag{1.1}
\]

The naive implementation materializes the $L \times L$ attention matrix, 
requiring $\mathcal{O}(L^2)$ memory. For $L = 8192$ and $d_h = 128$, this 
is 2 GB \textbf{per head, per layer} in FP16 — infeasible for most GPUs.

\textbf{FlashAttention-3} (Dao et al., 2024) solves this via \textbf{tiling} and 
\textbf{recomputation}: the attention matrix is computed in blocks that fit in 
GPU shared memory, and intermediate activations are recomputed during 
backward rather than stored. This gives:

\begin{itemize}
\item \textbf{Memory:} $\mathcal{O}(L)$ instead of $\mathcal{O}(L^2)$.
\end{itemize}
\begin{itemize}
\item \textbf{Throughput:} 2-4× faster than standard attention on Hopper GPUs.
\end{itemize}

\textbf{FP8 support.} FlashAttention-3 adds FP8 (e4m3) tensor core utilization, 
achieving another 2× on H100/H200.

\subsection{1.2 The Deployment Problem}

Despite its benefits, FlashAttention-3 is not a drop-in replacement:

1. \textbf{Hardware requirement.} Requires NVIDIA Hopper (H100/H200) for 
   TMA and FP8. Ampere (A100) supports FlashAttention-2 only; older GPUs 
   use the SDPA fallback.
2. \textbf{Software dependency.} Requires the \texttt{flash-attn} or \texttt{flash-attn-3} 
   Python package to be installed.
3. \textbf{Architecture heterogeneity.} Pretrained models differ in their 
   attention organization (fused Conv1D, nested self-attention, GQA, etc.), 
   requiring per-architecture weight extraction.

Our wrapper addresses all three:

\begin{itemize}
\item \textbf{Automatic fallback.} If FlashAttention-3 is not available, use
\end{itemize}
  PyTorch's SDPA kernel (which itself may use FlashAttention-2 internally).
\begin{itemize}
\item \textbf{Multi-architecture weight copy.} A single \texttt{copy_weights} handles
\end{itemize}
  GPT-2 Conv1D, RoBERTa/BERT nested, DistilBERT flat, and Qwen/LLaMA GQA.
\begin{itemize}
\item \textbf{FP8 opt-in.} FP8 precision is enabled via a flag, with automatic
\end{itemize}
  dtype conversion and graceful NaN handling.

\subsection{1.3 Problem Statement}

Given a pretrained attention module $\mathcal{A}$ with projections 
$W_Q, W_K, W_V, W_O$, we seek a wrapper $\mathcal{A}'$ such that:

1. \textbf{Exactness.} $\mathcal{A}'(X) = \mathcal{A}(X)$ mathematically (up to 
   numerical precision).
2. \textbf{Efficiency.} When FlashAttention-3 is available, 
   $\text{Time}(\mathcal{A}') < \text{Time}(\mathcal{A})$ on Hopper GPUs.
3. \textbf{Robustness.} When FlashAttention-3 is not available, gracefully fall 
   back to SDPA without crashing.
4. \textbf{Portability.} Supports flat, nested, and causal architectures, plus 
   GQA and RoPE.
5. \textbf{Recovery capacity.} Add trainable LoRA adapters for task adaptation.

\textbf{Note.} This is \textbf{not} a compression method. Parameter count increases 
by the LoRA adapter size.

\subsection{1.4 Our Approach}

We provide a wrapper class \texttt{FlashAttention3} that:

1. \textbf{Extracts pretrained weights} from the source attention module.
2. \textbf{Freezes base weights} and adds trainable LoRA adapters.
3. \textbf{Dispatches} to FlashAttention-3 if available, otherwise to SDPA.
4. \textbf{Supports GQA} via \texttt{repeat_kv} before the attention call.
5. \textbf{Applies RoPE} (with partial-dim support) before the attention call.
6. \textbf{Handles causal masking} for decoder models.
7. \textbf{Supports cross-attention} for encoder-decoder models.

\subsection{1.5 Contributions}

\begin{itemize}
\item \textbf{A production-hardened FlashAttention-3 wrapper.} Works across four
\end{itemize}
  architectures with exact weight transfer and graceful fallback.
\begin{itemize}
\item \textbf{Multi-architecture weight extraction.} Handles GPT-2 Conv1D,
\end{itemize}
  RoBERTa/BERT nested, DistilBERT flat, and Qwen/LLaMA GQA.
\begin{itemize}
\item \textbf{FP8 support with automatic fallback.} FP8 precision on Hopper GPUs
\end{itemize}
  with FP16 fallback if the hardware or library doesn't support it.
\begin{itemize}
\item \textbf{Honest reporting.} We disclose that on non-Hopper GPUs (like our
\end{itemize}
  RTX 4070 test machine), the SDPA fallback is used, not the FlashAttention-3 
  kernel.
\begin{itemize}
\item \textbf{Full reproducibility.} Code and evaluation scripts released under
\end{itemize}
  MIT.

\subsection{1.6 Paper Organization}

Section 2 reviews related work. Section 3 provides background on 
FlashAttention, SDPA, and hardware acceleration. Section 4 describes the 
wrapper. Section 5 details architecture integration. Section 6 gives 
experimental setup. Section 7 reports results. Section 8 analyzes the 
crossover between kernel and fallback paths. Section 9 discusses practical 
usage. Section 10 lists limitations. Section 11 concludes.



\subsection{Notation}
